\chapter{Chapter09 — Hypercontractive Necessity}
%%% generated by the blueprint pipeline from FABL.Chapter09.HypercontractiveNecessity
\section{Overview}

Declarations of the Lean module \texttt{FABL.Chapter09.HypercontractiveNecessity}.

\section{Declarations}

\begin{definition}[Affine moment]
\lean{FABL.affineMoment}
\statementsource{Analysis-of-Boolean-Functions-by-Ryan-ODonnell}{pdf-ffbe9c7b3306-p272-b002}
\label{FABL.affineMoment}
\leanok
Let $\Omega$ be a finite type, let $r,c,t\in\mathbb R$, and let $X:\Omega\to\mathbb R$.
Define the affine moment as the expectation over $x\in\Omega$ of $(1+ctX(x))^r$.
\end{definition}

\begin{definition}[Affine moment derivative expression]
\lean{FABL.affineMomentDeriv}
\statementsource{Analysis-of-Boolean-Functions-by-Ryan-ODonnell}{pdf-ffbe9c7b3306-p272-b002}
\label{FABL.affineMomentDeriv}
\leanok
Let $\Omega$ be a finite type, let $r,c,t\in\mathbb R$, and let $X:\Omega\to\mathbb R$.
Define the associated expression as the expectation over $x\in\Omega$ of
$rcX(x)(1+ctX(x))^{r-1}$.
\end{definition}

\begin{definition}[Second affine moment derivative expression]
\lean{FABL.affineMomentDeriv2}
\statementsource{Analysis-of-Boolean-Functions-by-Ryan-ODonnell}{pdf-ffbe9c7b3306-p272-b002}
\label{FABL.affineMomentDeriv2}
\leanok
Let $\Omega$ be a finite type, let $r,c,t\in\mathbb R$, and let $X:\Omega\to\mathbb R$.
Define the associated second-order expression as the expectation over $x\in\Omega$ of
$r(r-1)c^2X(x)^2(1+ctX(x))^{r-2}$.
\end{definition}

\begin{theorem}[Eventual positivity of a finite affine family]
\lean{FABL.eventually_affine_positive}
\label{FABL.eventually_affine_positive}
\leanok
\difficulty{4}
Let $\Omega$ be a finite set, let $c\in\mathbb R$, and let $X:\Omega\to\mathbb R$. For
all real $t$ sufficiently close to $0$, one has $1+ctX(x)>0$ for every $x\in\Omega$.
\end{theorem}

\begin{theorem}[Derivative of the affine moment]
\lean{FABL.affineMoment_hasDerivAt}
\label{FABL.affineMoment_hasDerivAt}
\leanok
\uses{FABL.affineMoment, FABL.affineMomentDeriv}
\difficulty{4}
Let $\Omega$ be a finite set, let $r,c\in\mathbb R$, and let $X:\Omega\to\mathbb R$.
Define $M(t)$ as the average over $x\in\Omega$ of $(1+ctX(x))^r$, and define $D(t)$ as
the average of $rcX(x)(1+ctX(x))^{r-1}$. At any $t\in\mathbb R$ for which $1+ctX(x)>0$
for every $x\in\Omega$, the function $M$ is differentiable at $t$ with derivative
$D(t)$.
\end{theorem}

\begin{theorem}[Derivative of the affine moment derivative]
\lean{FABL.affineMomentDeriv_hasDerivAt}
\label{FABL.affineMomentDeriv_hasDerivAt}
\leanok
\uses{FABL.affineMomentDeriv, FABL.affineMomentDeriv2}
\difficulty{4}
Let $\Omega$ be a finite set, let $r,c\in\mathbb R$, and let $X:\Omega\to\mathbb R$.
Define $F(t)$ as the average over $x\in\Omega$ of $rcX(x)(1+ctX(x))^{r-1}$, and define
$G(t)$ as the average of $r(r-1)c^2X(x)^2(1+ctX(x))^{r-2}$. At any $t\in\mathbb R$ for
which $1+ctX(x)>0$ for every $x\in\Omega$, the function $F$ is differentiable at $t$
with derivative $G(t)$.
\end{theorem}

\begin{theorem}[Affine moment at zero]
\lean{FABL.affineMoment_zero}
\label{FABL.affineMoment_zero}
\leanok
\uses{FABL.affineMoment}
\difficulty{1}
Let $\Omega$ be a finite nonempty set, let $r,c\in\mathbb R$, and let
$X:\Omega\to\mathbb R$. The average over $\Omega$ of $(1+ctX(x))^r$ at $t=0$ is $1$.
\end{theorem}

\begin{theorem}[Affine moment derivative at zero]
\lean{FABL.affineMomentDeriv_zero}
\label{FABL.affineMomentDeriv_zero}
\leanok
\uses{FABL.affineMomentDeriv}
\difficulty{3}
Let $\Omega$ be a finite nonempty set, let $r,c\in\mathbb R$, and let
$X:\Omega\to\mathbb R$. The affine moment derivative expression at $t=0$, defined as the
average of $rcX(x)(1+ctX(x))^{r-1}$, equals $rc$ times the average of $X(x)$.
\end{theorem}

\begin{theorem}[Second affine moment derivative at zero]
\lean{FABL.affineMomentDeriv2_zero}
\label{FABL.affineMomentDeriv2_zero}
\leanok
\uses{FABL.affineMomentDeriv2}
\difficulty{3}
Let $\Omega$ be a finite nonempty set, let $r,c\in\mathbb R$, and let
$X:\Omega\to\mathbb R$. The second derivative expression at $t=0$, defined as the
average over $x\in\Omega$ of $r(r-1)c^2X(x)^2(1+ctX(x))^{r-2}$, equals $r(r-1)c^2$ times
the average of $X(x)^2$.
\end{theorem}

\begin{definition}[Hypercontractive moment gap]
\lean{FABL.hypercontractiveMomentGap}
\statementsource{Analysis-of-Boolean-Functions-by-Ryan-ODonnell}{pdf-ffbe9c7b3306-p272-b002}
\label{FABL.hypercontractiveMomentGap}
\leanok
\uses{FABL.affineMoment}
For a finite set $\Omega$, a function $X\colon\Omega\to\mathbb{R}$, and real numbers
$p,q,\rho,t$, define the hypercontractive moment gap to be the difference between the
$q$th power of the affine moment of order $p$ with parameter $1$ and the $p$th power of
the affine moment of order $q$ with parameter $\rho$, both evaluated at $X$ and $t$.
\end{definition}

\begin{definition}[Hypercontractive moment gap derivative]
\lean{FABL.hypercontractiveMomentGapDeriv}
\statementsource{Analysis-of-Boolean-Functions-by-Ryan-ODonnell}{pdf-ffbe9c7b3306-p272-b002}
\label{FABL.hypercontractiveMomentGapDeriv}
\leanok
\uses{FABL.affineMoment, FABL.affineMomentDeriv}
For a finite set $\Omega$, a function $X:\Omega\to\mathbb{R}$, and real numbers
$p,q,\rho,t$, define the hypercontractive moment gap derivative to be
\[
q\,D_{p,1}(X,t)\,M_{p,1}(X,t)^{q-1}-p\,D_{q,\rho}(X,t)\,M_{q,\rho}(X,t)^{p-1},
\]
where $M_{a,r}(X,t)$ denotes the affine moment with parameters $a,r,X,t$, and
$D_{a,r}(X,t)$ its affine moment derivative.
\end{definition}

\begin{theorem}[Positivity of the affine moment]
\lean{FABL.affineMoment_pos}
\label{FABL.affineMoment_pos}
\leanok
\uses{FABL.affineMoment}
\difficulty{3}
Let $\Omega$ be a finite nonempty set, let $r,c,t\in\mathbb R$, and let
$X:\Omega\to\mathbb R$. If $1+ctX(x)>0$ for every $x\in\Omega$, then the average over
$\Omega$ of $(1+ctX(x))^r$ is positive.
\end{theorem}

\begin{theorem}[Derivative of the moment gap]
\lean{FABL.hypercontractiveMomentGap_hasDerivAt}
\label{FABL.hypercontractiveMomentGap_hasDerivAt}
\leanok
\uses{FABL.affineMoment_hasDerivAt, FABL.affineMoment_pos, FABL.hypercontractiveMomentGap, FABL.hypercontractiveMomentGapDeriv}
\difficulty{4}
Let $\Omega$ be a finite nonempty set, let $p,q,\rho,t\in\mathbb R$, and let
$X:\Omega\to\mathbb R$. Suppose that $1+tX(x)>0$ and $1+\rho tX(x)>0$ for every
$x\in\Omega$. Then the hypercontractive moment gap is differentiable at $t$, with
derivative equal to the difference between $q$ times the derivative of the first affine
moment raised to the power $q$, and $p$ times the derivative of the second affine moment
raised to the power $p$, as specified by the moment-gap derivative.
\end{theorem}

\begin{theorem}[Derivative of the moment gap at zero]
\lean{FABL.hypercontractiveMomentGapDeriv_hasDerivAt_zero}
\label{FABL.hypercontractiveMomentGapDeriv_hasDerivAt_zero}
\leanok
\uses{FABL.affineMomentDeriv2_zero, FABL.affineMomentDeriv_hasDerivAt, FABL.affineMomentDeriv_zero, FABL.affineMoment_hasDerivAt, FABL.affineMoment_zero, FABL.hypercontractiveMomentGapDeriv}
\difficulty{5}
Let $\Omega$ be a finite nonempty set, let $p,q,\rho\in\mathbb R$, and let
$X:\Omega\to\mathbb R$ have mean zero. The hypercontractive moment gap, as a function of
$t$, is differentiable at $t=0$, with derivative
\[pq\bigl((p-1)-(q-1)\rho^2\bigr)\,\mathbb E_{x\in\Omega}X(x)^2.\]
\end{theorem}

\begin{theorem}[Vanishing moment gap at zero]
\lean{FABL.hypercontractiveMomentGap_zero}
\label{FABL.hypercontractiveMomentGap_zero}
\leanok
\uses{FABL.hypercontractiveMomentGap}
\difficulty{1}
Let $\Omega$ be a finite nonempty set, let $p,q,\rho\in\mathbb{R}$, and let
$X\colon\Omega\to\mathbb{R}$. The difference between the $q$th power of the affine
moment of order $p$ with coefficient $1$ and the $p$th power of the affine moment of
order $q$ with coefficient $\rho$ is zero at $t=0$.
\end{theorem}

\begin{theorem}[Derivative of the moment gap at zero]
\lean{FABL.hypercontractiveMomentGap_hasDerivAt_zero}
\label{FABL.hypercontractiveMomentGap_hasDerivAt_zero}
\leanok
\uses{FABL.affineMomentDeriv_zero, FABL.affineMoment_zero, FABL.hypercontractiveMomentGap, FABL.hypercontractiveMomentGapDeriv, FABL.hypercontractiveMomentGap_hasDerivAt}
\difficulty{3}
Let $\Omega$ be a finite nonempty set, let $p,q,\rho\in\mathbb R$, and let
$X:\Omega\to\mathbb R$. The hypercontractive moment gap, as a function of $t$, is
differentiable at $t=0$, with derivative
\[pq(1-\rho)\,\mathbb E_{x\in\Omega}X(x).\]
\end{theorem}

\begin{theorem}[Affine moment as an absolute moment]
\lean{FABL.affineMoment_eq_absMoment}
\label{FABL.affineMoment_eq_absMoment}
\leanok
\uses{FABL.affineMoment}
\difficulty{3}
Let $\Omega$ be a finite set, let $r,c,t\in\mathbb R$, and let $X:\Omega\to\mathbb R$.
If $1+ctX(x)>0$ for every $x\in\Omega$, then the average of $(1+ctX(x))^r$ over $\Omega$
equals the average of $|1+ctX(x)|^r$.
\end{theorem}

\begin{theorem}[Hypercontractive moment gap nonneg]
\lean{FABL.hypercontractiveMomentGap_nonneg}
\statementsource{Analysis-of-Boolean-Functions-by-Ryan-ODonnell}{pdf-ffbe9c7b3306-p272-b002}
\label{FABL.hypercontractiveMomentGap_nonneg}
\leanok
\uses{FABL.IsHypercontractiveRandomVariable, FABL.affineMoment, FABL.affineMoment_eq_absMoment, FABL.hypercontractiveMomentGap, FABL.uniformLpNorm, FABL.uniformLpNormENNReal_coe, FABL.uniformLpNorm_nonneg, FABL.uniformLpNorm_rpow}
\difficulty{4}
% dataset-pipeline: stub -- prose pending build_informal_statements.py
Statement pending regeneration.
\end{theorem}

\begin{theorem}[Second derivative at a local minimum]
\lean{FABL.secondDeriv_nonneg_of_isLocalMin}
\label{FABL.secondDeriv_nonneg_of_isLocalMin}
\leanok
\difficulty{5}
Let $f\colon\mathbb{R}\to\mathbb{R}$ and let $a,d,d_2\in\mathbb{R}$. If $f$ has a local
minimum at $a$, its derivative at $a$ is $d$, and the derivative of $f'$ at $a$ is
$d_2$, then $d_2\ge 0$.
\end{theorem}

\begin{theorem}[Mean of a hypercontractive variable]
\lean{FABL.IsHypercontractiveRandomVariable.mean_eq_zero}
\proofsource{Analysis-of-Boolean-Functions-by-Ryan-ODonnell}{
  pdf-ffbe9c7b3306-p256-b003,
  pdf-ffbe9c7b3306-p256-b004,
  pdf-ffbe9c7b3306-p272-b002
}
\label{FABL.IsHypercontractiveRandomVariable.mean_eq_zero}
\leanok
\uses{FABL.IsHypercontractiveRandomVariable, FABL.eventually_affine_positive, FABL.hypercontractiveMomentGap, FABL.hypercontractiveMomentGap_hasDerivAt_zero, FABL.hypercontractiveMomentGap_nonneg}
\difficulty{4}
Let $\\Omega$ be a finite nonempty set, let $p,q>0$ be real numbers, let $\\rho<1$, and
let $X:\\Omega\\to\\mathbb R$. Suppose that for every $a,b\\in\\mathbb R$, the uniform
$L^q$ norm of $a+\\rho bX$ is at most the uniform $L^p$ norm of $a+bX$. Then the mean of
$X$ on $\\Omega$ is zero. Here the uniform norms are taken with respect to the uniform
distribution on $\\Omega$.
\end{theorem}

\begin{theorem}[Squared contraction parameter bound]
\lean{FABL.IsHypercontractiveRandomVariable.rho_sq_bound}
\proofsource{Analysis-of-Boolean-Functions-by-Ryan-ODonnell}{pdf-ffbe9c7b3306-p272-b002}
\label{FABL.IsHypercontractiveRandomVariable.rho_sq_bound}
\leanok
\uses{FABL.IsHypercontractiveRandomVariable, FABL.IsHypercontractiveRandomVariable.mean_eq_zero, FABL.eventually_affine_positive, FABL.hypercontractiveMomentGap, FABL.hypercontractiveMomentGapDeriv, FABL.hypercontractiveMomentGapDeriv_hasDerivAt_zero, FABL.hypercontractiveMomentGap_hasDerivAt, FABL.hypercontractiveMomentGap_hasDerivAt_zero, FABL.hypercontractiveMomentGap_nonneg, FABL.secondDeriv_nonneg_of_isLocalMin}
\difficulty{6}
Let $\Omega$ be a finite nonempty set, let $p,q\ge0$ satisfy $p>0$ and $q>0$, let
$\rho<1$, and let $X:\Omega\to\mathbb R$. Suppose the uniform $q$-norm of $x\mapsto
a+\rho bX(x)$ is at most the uniform $p$-norm of $x\mapsto a+bX(x)$ for all real $a,b$,
and suppose the uniform expectation of $X^2$ is positive. Then $(q-1)\rho^2\le p-1$.
\end{theorem}

\begin{theorem}[Contraction parameter bounded by a square-root exponent ratio]
\lean{FABL.IsHypercontractiveRandomVariable.rho_le_sqrt_ratio}
\proofsource{Analysis-of-Boolean-Functions-by-Ryan-ODonnell}{pdf-ffbe9c7b3306-p272-b002}
\label{FABL.IsHypercontractiveRandomVariable.rho_le_sqrt_ratio}
\leanok
\uses{FABL.IsHypercontractiveRandomVariable, FABL.IsHypercontractiveRandomVariable.rho_sq_bound}
\difficulty{3}
Let $\Omega$ be a finite nonempty set, let $p,q\ge0$ satisfy $p\ge1$ and $q>1$, let
$\rho<1$, and let $X:\Omega\to\mathbb R$. Suppose the uniform $q$-norm of $x\mapsto
a+\rho bX(x)$ is at most the uniform $p$-norm of $x\mapsto a+bX(x)$ for all real $a,b$,
and suppose the uniform expectation of $X^2$ is positive. Then
$\rho\le\sqrt{(p-1)/(q-1)}$.
\end{theorem}

\begin{theorem}[Hypercontractivity of the Zero Variable]
\lean{FABL.zeroVariable_one_two_half_hypercontractive}
\label{FABL.zeroVariable_one_two_half_hypercontractive}
\leanok
\uses{FABL.IsHypercontractiveRandomVariable, FABL.uniformLpNorm, FABL.uniformLpNormENNReal, FABL.uniformLpNormENNReal_coe, FABL.uniformLpNorm_const_mul}
\difficulty{4}
For all real numbers $a$ and $b$, the uniform $L^2$ norm on a one-point space of the
function with constant value $a+\tfrac12 b\cdot 0$ is at most the uniform $L^1$ norm of
the function with constant value $a+b\cdot 0$.
\end{theorem}

\begin{theorem}[Failure at a zero numerator]
\lean{FABL.exercise9_10_nonzero_hypothesis_necessary}
\label{FABL.exercise9_10_nonzero_hypothesis_necessary}
\leanok
\difficulty{2}
The inequality $\frac12\le\sqrt{(1-1)/(2-1)}$ is false.
\end{theorem}

\begin{theorem}[Positive second moment from positive norm]
\lean{FABL.expect_sq_pos_of_uniformLpNorm_two_pos}
\label{FABL.expect_sq_pos_of_uniformLpNorm_two_pos}
\leanok
\uses{FABL.uniformLpNorm, FABL.uniformLpNorm_two_sq_eq_expect_sq}
\difficulty{3}
Let $\Omega$ be a finite nonempty set and let $X\colon\Omega\to\mathbb{R}$. If the $L^2$
norm of $X$ under the uniform probability measure on $\Omega$ is positive, then its
second moment is positive: $\mathbb{E}_{x\in\Omega}[X(x)^2]>0$.
\end{theorem}

\begin{theorem}[Anticoncentration under hypercontractivity]
\lean{FABL.translated_anticoncentration_exact_of_hypercontractive}
\proofsource{Analysis-of-Boolean-Functions-by-Ryan-ODonnell}{pdf-ffbe9c7b3306-p272-b004}
\label{FABL.translated_anticoncentration_exact_of_hypercontractive}
\leanok
\uses{FABL.IsHypercontractiveRandomVariable, FABL.IsHypercontractiveRandomVariable.mean_eq_zero, FABL.translated_anticoncentration_exact, FABL.uniformLpNorm, FABL.uniformProbability}
\difficulty{3}
Let $\Omega$ be a finite nonempty set, and let $X:\Omega\to\mathbb R$ have positive
uniform $L^2$ norm $M$. Suppose $0<\rho<1$ and, for every $a,b\in\mathbb R$, the uniform
$L^4$ norm of $a+\rho bX$ is at most the uniform $L^2$ norm of $a+bX$. Then for every
$\theta\in\mathbb R$ and every $0<t<1$,
\[
\left(\frac{\rho^2(1-t^2)+\rho^2(\theta/M)^2}{1+\rho^2(\theta/M)^2}\right)^2
\leq \mathbb P_{x\text{ uniform on }\Omega}\!\left(tM<|X(x)-\theta|\right).
\]
\end{theorem}

\begin{theorem}[Anticoncentration under hypercontractivity]
\lean{FABL.translated_anticoncentration_of_hypercontractive}
\proofsource{Analysis-of-Boolean-Functions-by-Ryan-ODonnell}{pdf-ffbe9c7b3306-p272-b004}
\label{FABL.translated_anticoncentration_of_hypercontractive}
\leanok
\uses{FABL.IsHypercontractiveRandomVariable, FABL.IsHypercontractiveRandomVariable.mean_eq_zero, FABL.IsHypercontractiveRandomVariable.rho_le_sqrt_ratio, FABL.expect_sq_pos_of_uniformLpNorm_two_pos, FABL.translated_anticoncentration, FABL.uniformLpNorm, FABL.uniformProbability}
\difficulty{4}
Let $\Omega$ be a finite nonempty set, and let $X:\Omega\to\mathbb R$ have positive
uniform $L^2$ norm. Suppose $0<\rho<1$ and, for every $a,b\in\mathbb R$, the uniform
$L^4$ norm of $x\mapsto a+\rho bX(x)$ is at most the uniform $L^2$ norm of $x\mapsto
a+bX(x)$. For $\theta\in\mathbb R$ and $0<t<1$, the uniform probability that
$|X(x)-\theta|>t\,\|X\|_2$ is at least $(1-t^2)^2\rho^4$, where $\|X\|_2$ denotes the
uniform $L^2$ norm.
\end{theorem}
